\begin{abstract}
% claims: C01 C03 C08 C09 C10 C11; sources: S03 S04 S10 S14 S15 S16 S17
Receipts can attest that trusted infrastructure recorded an internal state, while
text watermarks can place a detectable statistical signal inside generated
output. Neither mechanism alone shows that the signal corresponds to a
causally operative state in the computation that produced the output. We ask
whether an authenticated internal state can select such a signal when the
prompt, final answer, surface content, generation settings, and underlying
random draws are held fixed. We study this question in an engineered arithmetic
model whose answer must pass through a supervised categorical bottleneck.
Trusted instrumentation authenticates the bottleneck state, derives a
state-specific carrier identity, and uses that identity to bias sampling in
grammar-constrained text. Our decisive comparison pairs a natural execution
with an intervention that changes the authenticated state and downstream
computation while preserving the answer. The public receipt-bound integration
was prospectively validated on 128/128 held-out public closures. Protected
qualification was not earned: one of 128 closures failed computation fidelity,
the descriptive result was 127/128, and final protected confirmation was not
run. This demonstrates bounded feasibility of state-conditioned statistical
output carriage in one selected engineered system, not output-only
authentication, whole-reasoning provenance, or protected generalisation.
\end{abstract}
